\section{问题分析}

\subsection{问题一的分析}
问题一研究前 $1800\,\mathrm{s}$ 的预热阶段。附件1给出离散环境数据，题面要求输出指定时刻和径向位置的温度与水分浓度。因此，本问的核心是将环境序列转化为连续边界，在固定半径的一维径向圆柱模型中分别求解导热和水分扩散过程，并以守恒离散保证表面通量与内部变化相容。

\subsection{问题二的分析}
问题二将时间范围扩展至前三小时，并采用附录3的变物性关系。相较问题一，密度、比热容和导热系数随含水率变化，水分扩散系数同时依赖含水率和温度，温度场与水分场由此形成参数耦合。本问仍采用固定半径和相同环境边界，重点分析温度快速响应与水分扩散滞后的差异。

\subsection{问题三的分析}
问题三要求确定全域水分浓度达到阈值的时间。它继承问题二的固定域耦合模型，仅增加长时段积分和全域事件判定。判停量取轴线、控制体及真实表面水分浓度的最大值，避免以表面或单一采样点代替整个药材的达标条件；附件1结束后的环境量按末值延拓，并对该边界假设进行敏感性分析。

\subsection{问题四的分析}
问题四同时改变几何和物性条件。附件2给出半径随时间的变化，附录4给出新的经验物性；固定物理位置还可能随收缩移出药材域。本问用材料坐标 $\xi=r/R(t)$ 将移动区域映射到固定区间，使收缩体现为传递尺度 $R(t)^{-2}$，并在输出时区分固定位置和实时表面。除比较最终临界时间外，还需检验移动模型在固定域和零扩散极限下的自洽性。

\subsection{统一求解与结果组织}
四问采用统一的轴对称径向传递框架和守恒有限体积离散；问题二至四同步更新状态相关物性，问题三、四以全域最大水分浓度触发连续事件。结果章节按“模型的建立—模型的求解—结果的得出与分析”组织，数值可靠性和敏感性在统一章节中汇总。

\begin{table}[H]
\centering
\caption{四个问题的建模递进与主要输出}
\label{tab:problem-analysis}
\zihao{-5}
\begin{tabular}{@{}lccc@{}}
\toprule
问题 & 几何与物性 & 主要计算量 & 主要输出 \\
\midrule
一 & 固定半径、附录2 & 解耦温度与水分场 & 指定时刻和位置场值 \\
二 & 固定半径、附录3 & 变物性耦合场 & 前三小时场值 \\
三 & 固定半径、附录3 & 全域事件积分 & 达标时间与径向分布 \\
四 & 径向收缩、附录4 & 移动域耦合场 & 达标时间、表面及影响对照 \\
\bottomrule
\end{tabular}
\end{table}
